% !TEX root = thesis-main.tex
\documentclass[12pt,a4paper,openright,twoside]{book}
\usepackage[utf8]{inputenc}
\usepackage{disi-thesis}
\usepackage{code-lstlistings}
\usepackage{notes}
\usepackage{shortcuts}
\usepackage{acronym}

\school{\unibo}
\programme{Corso di Laurea in Ingegneria e Scienze Informatiche}
\title{Fancy Title}
\author{Candidate Name}
\date{\today}
\subject{Nome del corso del relatore}
\supervisor{Prof. Supervisor Here}
% \cosupervisor{Dott. CoSupervisor 1}       % uncomment if needed
% \morecosupervisor{Dott. CoSupervisor 2}   % uncomment if needed
\session{I}
\academicyear{2024-2025}

% Definition of acronyms
\acrodef{fmi}[FMI]{Functional Mock-up Interface}
\acrodef{fmu}[FMU]{Functional Mock-up Unit}
\acrodef{kmp}[KMP]{Kotlin Multiplatform}
\acrodef{api}[API]{Application Programming Interface}
\acrodef{ci}[CI]{Continuous Integration}
\acrodef{cd}[CD]{Continuous Delivery}
\acrodef{rest}[REST]{Representational State Transfer}
\acrodef{ui}[UI]{User Interface}
\acrodef{me}[ME]{Model Exchange}
\acrodef{cs}[CS]{Co-Simulation}
\acrodef{jvm}[JVM]{Java Virtual Machine}
\acrodef{cio}[CIO]{Coroutine-based I/O}

\mainlinespacing{1.241}

\begin{document}

\frontmatter\frontispiece

\begin{abstract}
% Max 2000 caratteri. Struttura consigliata: Contesto, Problema/Obiettivi, Metodi/Contributo, Risultati, Conclusioni.
\end{abstract}

% \begin{dedication}
% Opzionale. Poche righe al massimo.
% \end{dedication}

%----------------------------------------------------------------------------------------
\tableofcontents
\listoffigures        % (optional) commentare se vuota
\lstlistoflistings    % (optional) commentare se vuota
%----------------------------------------------------------------------------------------

\mainmatter

%----------------------------------------------------------------------------------------
\chapter{Introduzione}
\label{chap:introduction}
%----------------------------------------------------------------------------------------

\section{Contesto e motivazione}
\label{sec:intro-contesto}

\section{Obiettivi del progetto}
\label{sec:intro-obiettivi}

\section{Struttura della tesi}
\label{sec:intro-struttura}

%----------------------------------------------------------------------------------------
\chapter{Background}
\label{chap:background}
%----------------------------------------------------------------------------------------

\section{Lo standard FMI e il formato FMU}
\label{sec:background-fmi}

\subsection{Functional Mock-up Interface (FMI 2.0)}
\label{subsec:background-fmi-standard}

\subsection{Struttura di un file FMU e modalità operative}
\label{subsec:background-fmu-struttura}

\section{Kotlin Multiplatform}
\label{sec:background-kmp}

\subsection{Architettura expect/actual e compilazione nativa}
\label{subsec:background-kmp-expect-actual}

\subsection{Interoperabilità con C tramite cinterop}
\label{subsec:background-kmp-cinterop}

\section{Tecnologie adottate}
\label{sec:background-tecnologie}

\subsection{FMILib}
\label{subsec:background-fmilib}

\subsection{Ktor}
\label{subsec:background-ktor}

\subsection{Compose Multiplatform}
\label{subsec:background-compose}

\section{Strumenti di processo}
\label{sec:background-strumenti}

\subsection{Pipeline CI/CD con GitHub Actions}
\label{subsec:background-cicd}

\subsection{Gradle e gestione multimodulo}
\label{subsec:background-gradle}

%----------------------------------------------------------------------------------------
\chapter{Analisi}
\label{chap:analisi}
%----------------------------------------------------------------------------------------

\section{Obiettivi}
\label{sec:analisi-obiettivi}

\section{Requisiti funzionali}
\label{sec:analisi-funzionali}

\section{Requisiti non funzionali}
\label{sec:analisi-nonfunzionali}

\section{Modello del dominio}
\label{sec:analisi-dominio}

%----------------------------------------------------------------------------------------
\chapter{Design}
\label{chap:design}
%----------------------------------------------------------------------------------------

\section{Architettura generale e struttura}
\label{sec:design-architettura}

\section{Modulo \texttt{fmilib}: integrazione della libreria C nativa}
\label{sec:design-fmilib}

\section{Modulo \texttt{fmu-kt}}
\label{sec:design-fmukt}

\subsection{Astrazione \texttt{FmuPreprocessor} e factory multipiattaforma}
\label{subsec:design-preprocessor}

\subsection{\texttt{NativeFmiWrapper}: ciclo di vita, metadati e simulazione}
\label{subsec:design-wrapper}

\subsection{\texttt{DefaultFmu} come API di alto livello}
\label{subsec:design-defaultfmu}

\section{Modulo \texttt{backend}: architettura del server e routing REST}
\label{sec:design-backend}

\section{Modulo \texttt{frontend}: struttura della UI e comunicazione con il backend}
\label{sec:design-frontend}

%----------------------------------------------------------------------------------------
\chapter{Implementazione}
\label{chap:implementazione}
%----------------------------------------------------------------------------------------

\section{Binding C/Kotlin e configurazione del linker per piattaforma}
\label{sec:impl-binding}

\section{Gestione del ciclo di vita dell'FMU}
\label{sec:impl-lifecycle}

\section{Estrazione dei metadati e gestione delle variabili}
\label{sec:impl-metadati}

\section{Esecuzione della simulazione}
\label{sec:impl-simulazione}

\section{Ricompilazione degli FMU su macOS ARM64}
\label{sec:impl-macos}

%----------------------------------------------------------------------------------------
\chapter{Valutazione}
\label{chap:valutazione}
%----------------------------------------------------------------------------------------

\section{Test unitari e di integrazione}
\label{sec:val-test}

\section{Copertura multipiattaforma nella CI}
\label{sec:val-ci}

%----------------------------------------------------------------------------------------
\chapter{Conclusioni e Sviluppi Futuri}
\label{chap:conclusioni}
%----------------------------------------------------------------------------------------

\section{Obiettivi raggiunti}
\label{sec:conclusioni-obiettivi}


\section{Sviluppi futuri}
\label{sec:conclusioni-futuro}

%----------------------------------------------------------------------------------------
% BIBLIOGRAPHY
%----------------------------------------------------------------------------------------
\backmatter
\nocite{*} % Rimuovere appena si ha la prima citazione
\bibliographystyle{alpha}
\bibliography{bibliography}

% \begin{acknowledgements}
% Opzionale. Max 1 pagina.
% \end{acknowledgements}

\end{document}